\section{Background}
\label{sec:background}

In this section we briefly describe the transformer architecture~\cite{Vaswani2017AttentionIA}.
We describe our modifications enabling set prediction in sec.~\ref{sec:method}, to which a reader familiar with
transformer architecture can skip directly.

\textit{Multi-head attention} is the core operation of the transformer.
We apply scaled dot-product multi-head attention, defined as:
\[\text{Attention}(\mQ, \mK, \mV) = \text{softmax}(\frac{\mQ\mK^T}{\sqrt{d}})\mV,\]
where $\mQ$, $\mK$ and $\mV$ are input matrices called query, key and value respectively,
and $d$ is number of channels in $\mK$.
In self-attention query, key and value matrices are obtained by projecting
the same input $\mX$ with different weights $W_{q,k,v}$:
$\mQ = \mW_q\mX$, $\mK = \mW_k\mX$, $\mV = \mW_v\mX$.
In inter-attention query is a different input.
Additionally, the multi-head aspect is achieved by doing attention multiple times with different weights on the same
inputs, and concatenating the outputs.

Let's define self-attention for 2D inputs, corresponding to activations of a convolutional network.
In this case input $\mX$ has shape $NHWC$, where $N$ is batch size,
$HW$ is activation height and width, and $C$ is the number of channels.
Query, key and value is then the same matrix, thus the name of the function.
An \textit{attention map} is a product of $\mX$ projected into $\mQ$ and $\mK$ after softmax,
computed for each sample in a minibatch and has $HW\times HW$ shape.
Self-attention thus has ability to `attend' to specific regions in the image, given the whole scene,
and utilize context.
Current region-based detectors lack this ability due to local region-detection relationship assumption.
In decoder inter-attention the size of the attention map computed is $H\times W$.

We adopt a common encoder-decoder architecture.
Encoder consists of multiple multi-head self-attention layers followed by a feed-forward network (FFN, or MLP) and layer normalization:
\begin{align}
    \mQ_\text{MHA} &= \text{layernorm}\Big(\mQ + \text{dropout}\big(\text{Attention}(\mQ, \mK, \mV)
        \big)\Big),\\
    \mQ &= \text{layernorm}(\mQ_\text{MHA} + \text{dropout}(\mW_2\text{relu}(\mW_1\mQ_\text{MHA}))),
    \label{eq:mha}
\end{align}
$\mW_1$ and $\mW_2$ are FFN weights.
We pass CNN activations through a single $1\times1$ convolutional layer and flatten
before feeding it to encoder as $\mQ$, $\mK$ and $\mV$.
Self-attention in encoder allows learning the regions of interest for classifying and localizing the object given the whole image.
Encoder output is called memory, which is fed into decoder together with $N$ query embeddings.
For each query embedding decoder outputs an object, or none, in parallel, as described later.
Decoder constists of multiple self-attention layers on query outputs,
inter-attention between query outputs and encoder memory, followed by FFN.
In inter-attention, $\mK$ and $\mQ$ are query outputs, and $\mV$ is encoder memory.